% Zone „Das kennst du schon“ – aus bank/kurvenuntersuchung/zone.jsonl (zusammenbau v0.1)
\einheitenkopf[zone][Kennst du schon]{Das kennst du schon}
\setcounter{aufgabe}{0}

%% TODO Titel: Ich-kann-Satz fehlt in der Bank (Platzhalter: Kettenname) – Nr. 1
%% TODO Anweisung: ein Satz über der ersten Teilaufgabe fehlt in der Bank – Nr. 1
\begin{aufgabe}{Ableitungen bilden}
\begin{teile}
\teil Bilde die Ableitung von $f(x) = x^5$. $f'(x) = $ \leerfeld
\teil Bilde die Ableitung von $f(x) = \frac{1}{4}x^4 - 2x^3 + 5$. $f'(x) = $ \leerfeld
\end{teile}
\end{aufgabe}

%% TODO Titel: Ich-kann-Satz fehlt in der Bank (Platzhalter: Kettenname) – Nr. 2
%% TODO Anweisung: ein Satz über der ersten Teilaufgabe fehlt in der Bank – Nr. 2
\begin{aufgabe}{Gleichungen lösen}
\begin{gleichungsraster}[2]
\gl{\text{Löse $3x - 12 = 0$.}} &
\gl{\text{Löse $x^2 - x - 0{,}75 = 0$ mit der Lösungsformel.}} \\
\end{gleichungsraster}
\end{aufgabe}

%% TODO Titel: Ich-kann-Satz fehlt in der Bank (Platzhalter: Kettenname) – Nr. 3
%% TODO Anweisung: ein Satz über der ersten Teilaufgabe fehlt in der Bank – Nr. 3
\begin{aufgabe}{Charakteristische Punkte (Achsenschnitt-, Hoch-, Tief-, Wendepunkte) am Graphen ganzrationaler Funktionen ablesen und im Sachzusammenhang deuten}
\begin{teile}
\teil Die Abbildung zeigt den Graphen von $f$. Lies die Koordinaten des Hochpunkts ab.

\begin{ksys}[xmin=-3,xmax=3,ymin=-3,ymax=3,ablesen] \funktion{-0.5*\x^3+1.5*\x+1}{f} \end{ksys}
$H($\leerfeld$|$\leerfeld$)$
\teil Die Abbildung zeigt den Graphen von $f$. Lies den Wendepunkt ab – dort, wo der Graph zwischen Tief- und Hochpunkt am steilsten ist.

\begin{ksys}[xmin=-3,xmax=3,ymin=-3,ymax=3,ablesen] \funktion{-0.5*\x^3+1.5*\x+1}{f} \end{ksys}
$W($\leerfeld$|$\leerfeld$)$
\end{teile}
\end{aufgabe}

%% TODO Titel: Ich-kann-Satz fehlt in der Bank (Platzhalter: Kettenname) – Nr. 4
%% TODO Anweisung: ein Satz über der ersten Teilaufgabe fehlt in der Bank – Nr. 4
\begin{aufgabe}{Funktionswerte berechnen und Punktprobe}
\begin{teile}
\teil Berechne $f(2)$ für $f(x) = x^2 + 3$. $f(2) = $ \leerfeld
\teil Berechne $f(-2)$ für $f(x) = 0{,}5x^3 - x$. $f(-2) = $ \leerfeld
\end{teile}
\end{aufgabe}

%% TODO Titel: Ich-kann-Satz fehlt in der Bank (Platzhalter: Kettenname) – Nr. 5
%% TODO Anweisung: ein Satz über der ersten Teilaufgabe fehlt in der Bank – Nr. 5
\begin{aufgabe}{Verlauf quadratischer und kubischer Graphen grob skizzieren (Öffnung, Symmetrie, Grenzverhalten)}
\begin{teile}
\teil Kreuze an: Die Parabel zu $f(x) = -2x^2 + 3$ ist \\ \kreuz{nach oben geöffnet} \\ \kreuz{nach unten geöffnet}
\teil Beschreibe, wie der Graph von $f(x) = -0{,}5x^3 + 2x$ für $x \to +\infty$ und für $x \to -\infty$ verläuft.
\end{teile}
\end{aufgabe}

%% TODO Titel: Ich-kann-Satz fehlt in der Bank (Platzhalter: Kettenname) – Nr. 6
%% TODO Anweisung: ein Satz über der ersten Teilaufgabe fehlt in der Bank – Nr. 6
\begin{aufgabe}{Ableitung als Tangentensteigung und lokale Änderungsrate deuten (Grundvorstellung)}
\begin{teile}
\teil Die Abbildung zeigt den Graphen von $f$. Ist die Steigung des Graphen im Punkt $A$ positiv, negativ oder null? \\ \kreuz{positiv} \kreuz{negativ} \kreuz{null}

\begin{ksys}[xmin=-3,xmax=4,ymin=-2,ymax=4,ablesen] \funktion{0.25*\x^2-1}{f} \punkt{2}{0}{A} \punkt{0}{-1}{B} \end{ksys}
\teil Die Abbildung zeigt den Graphen von $f$ mit der Tangente an der Stelle $2$. Lies die Steigung der Tangente ab: Wie groß ist $f'(2)$?

\begin{ksys}[xmin=-1,xmax=4,ymin=-1,ymax=5,ablesen] \funktion{0.5*\x^2}{f} \tangentean{0.5*\x^2}{2}{t} \end{ksys}
$f'(2) = $ \leerfeld
\end{teile}
\end{aufgabe}

%% TODO Titel: Ich-kann-Satz fehlt in der Bank (Platzhalter: Kettenname) – Nr. 7
%% TODO Anweisung: ein Satz über der ersten Teilaufgabe fehlt in der Bank – Nr. 7
\begin{aufgabe}{Gleichungen lösen – Fehler finden}
\begin{teile}
\teil Finde den Fehler und rechne richtig. \rechnung{x^3 - 5x^2 &= 0 &&\mid : x^2 \\ x - 5 &= 0 \\ x &= 5}
\end{teile}
\end{aufgabe}

%% TODO Titel: Ich-kann-Satz fehlt in der Bank (Platzhalter: Kettenname) – Nr. 8
%% TODO Anweisung: ein Satz über der ersten Teilaufgabe fehlt in der Bank – Nr. 8
\begin{aufgabe}{Gleichungen lösen – selbst rechnen}
\begin{gleichungsraster}[2]
\gl{\text{Löse $x^4 - 6x^3 = 0$.}} & \\
\end{gleichungsraster}
\end{aufgabe}
